\chapter{Algebraic convergence}

\section{Limits of discrete groups}

\begin{definition}[Geometric convergence of groups]
\label{thurston-ch09-geometric-convergence}
\group{thurston.ch09.convergence}\source{thurston:page-0241}\notready
 A sequence of closed subgroups $\Gamma_i$ of a Lie group $G$ converges geometrically to $\Gamma$ if every $\gamma\in\Gamma$ is the limit of elements $\gamma_i\in\Gamma_i$, and every limit of a convergent sequence $\gamma_{i_j}\in\Gamma_{i_j}$ belongs to $\Gamma$. The resulting topology on closed subgroups is the geometric topology.
\end{definition}

\begin{definition}[Algebraic and strong convergence]
\label{thurston-ch09-algebraic-strong-convergence}
\group{thurston.ch09.convergence}\source{thurston:page-0241}\source{thurston:page-0242}\notready
 For a countable group $\Gamma$ and representations $\rho_i:\Gamma\to G$, algebraic convergence means pointwise convergence $\rho_i(\gamma)\to\rho(\gamma)$ for every $\gamma$. Strong convergence means algebraic convergence together with geometric convergence of the image groups to $\rho(\Gamma)$.
\end{definition}

\begin{proposition}[Algebraic limits lie in geometric limits]
\label{thurston-ch09-algebraic-in-geometric}
\dcref{9.1.5}\group{thurston.ch09.convergence}\source{thurston:page-0244}\notready
 If $\rho_i\to\rho$ algebraically and $\rho_i(\Gamma)\to\Gamma'$ geometrically, then $\rho(\Gamma)\subseteq\Gamma'$.
\end{proposition}
\begin{proof}\sketch Pointwise convergence supplies, for each element of $\rho(\Gamma)$, a convergent sequence of elements of the image groups; apply the first clause of geometric convergence.\end{proof}

\begin{proposition}[Compactness of closed subgroups]
\label{thurston-ch09-closed-subgroup-compactness}
\dcref{9.1.6}\group{thurston.ch09.convergence}\source{thurston:page-0244}\notready
 For every Lie group $G$, its space of closed subgroups is compact in the geometric topology.
\end{proposition}
\begin{proof}\sketch If small nonidentity elements have a uniform lower bound, diagonal compactness of the bounded finite sets gives a discrete limit. Otherwise select the maximal limiting subspace of infinitesimal directions; it is closed under Lie brackets, and the same diagonal argument gives a closed subgroup with that tangent space.\end{proof}

\begin{corollary}[Compactness of based thick hyperbolic manifolds]
\label{thurston-ch09-based-thick-compactness}
\dcref{9.1.7}\group{thurston.ch09.convergence}\source{thurston:page-0244}\notready
 The space of complete hyperbolic manifolds equipped with a base frame in the $\epsilon$-thick part is compact in the geometric topology.
\end{corollary}

\begin{corollary}[Discreteness of non-elementary algebraic limits]
\label{thurston-ch09-non-elementary-limit-discrete}
\dcref{9.1.8}\group{thurston.ch09.convergence}\source{thurston:page-0245}\notready
 Let $\Gamma$ be countable and let discrete representations $\rho_i:\Gamma\to\PSL(2,\C)$ converge algebraically to $\rho$. If $\rho(\Gamma)$ has no abelian subgroup of finite index, then a subsequence converges geometrically to a discrete group containing $\rho(\Gamma)$; in particular $\rho(\Gamma)$ is discrete.
\end{corollary}
\begin{proof}\sketch Use compactness of based thick manifolds to obtain a closed geometric limit. Its identity component is abelian by the small-subgroup theorem, hence is trivial because it contains the non-elementary algebraic limit.\end{proof}

\begin{theorem}[Strong convergence of quasi-Fuchsian limits]
\label{thurston-ch09-quasi-fuchsian-strong-limit}
\dcref{9.2}\group{thurston.ch09.convergence}\source{thurston:page-0245}\source{thurston:page-0246}\notready
 Let $S$ have finite area and let faithful quasi-Fuchsian representations $\rho_i:\pi_1S\to\PSL(2,\C)$ preserve parabolicity and converge algebraically to $\rho$, with no new parabolic elements. Then $N_{\rho(\pi_1S)}$ is geometrically tame and the convergence is strong.
\end{theorem}
\begin{proof}\sketch Compact bending data either gives two disjoint locally convex limiting boundary surfaces, proving geometric finiteness and then strong convergence, or escapes compact sets, producing a geometrically infinite tame end. The covering restriction below prevents both ends from collapsing into one.\end{proof}

\begin{definition}[Almost geometrically tame end]
\label{thurston-ch09-almost-geometrically-tame}
\group{thurston.ch09.ends}\source{thurston:page-0246}\notready
 For a cusped hyperbolic manifold $N$, an end of $N-P$ is almost geometrically tame if a finite-sheeted cover of it agrees, outside a compact set, with a geometrically tame end.
\end{definition}

\begin{theorem}[Covering restriction for tame ends]
\label{thurston-ch09-covering-restriction}
\dcref{9.2.2}\group{thurston.ch09.ends}\source{thurston:page-0246}\source{thurston:page-0247}\notready
 Let $\widetilde N\to N$ be a cover and let $E$ be a geometrically infinite tame end of $\widetilde N-\widetilde P$, bounded by $S_{[e,\infty)}$. Either $N$ has finite volume and a finite cover fibers over $S^1$ with fiber $S$, or the image of $E$ in $N-P$, modulo a compact set, is almost geometrically tame.
\end{theorem}
\begin{proof}\sketch If the projection has finite local degree, covering transformations preserving the end contain $\pi_1S$ with finite index. Otherwise intersecting uncrumpled surfaces have precompact images; noncompactness yields a fibered finite cover, while compactness forces infinitely many compact identifications and reduces to the finite-degree case.\end{proof}

\section{The ending of an end}

\begin{theorem}[Ruled-cylinder area estimate]
\label{thurston-ch09-ruled-cylinder-area}
\dcref{9.3.1}\group{thurston.ch09.ending}\source{thurston:page-0249}\notready
 A ruled cylinder spanning homotopic closed geodesics has area bounded above by the sum of the lengths of its boundary curves.
\end{theorem}
\begin{proof}\sketch Rule the cylinder by geodesic segments and compare the induced Jacobi-field lengths; integrating the convexity estimate gives the boundary-length bound.\end{proof}

\begin{theorem}[Truncated ruled-cylinder estimate]
\label{thurston-ch09-truncated-cylinder-area}
\dcref{9.3.2}\group{thurston.ch09.ending}\source{thurston:page-0250}\notready
 If $C$ is the ruled cylinder between geodesics $\gamma_1,\gamma_2$, then for $r\ge0$,
\[
 \Area(C-\mathcal N_r\gamma_1)\le e^{-r}\length(\gamma_1)+\length(\gamma_2).
\]
\end{theorem}
\begin{proof}\sketch The contribution of the ruled sides decays exponentially with distance from $\gamma_1$, while the opposite boundary contributes at most its length.\end{proof}

\begin{proposition}[Intersection estimate in an end]
\label{thurston-ch09-end-intersection-estimate}
\dcref{9.3.3}\group{thurston.ch09.ending}\source{thurston:page-0252}\notready
 For curves $\alpha,\beta$ in an end and their associated cylinders, $|\alpha\cap C_\beta|+|C_\alpha\cap\beta|\ge i(\alpha',\beta')$, where $\alpha',\beta'$ are the corresponding boundary curves.
\end{proposition}

\begin{proposition}[Curve intersection estimate in an end]
\label{thurston-ch09-short-curve-bound}
\dcref{9.3.4}\group{thurston.ch09.ending}\source{thurston:page-0252,thurston:page-0253}\notready
 Let $E$ be a geometrically tame end cut off by $S_{[\epsilon,\infty)}$.
There is a constant $K$ such that the following holds.  Suppose that a curve
$\alpha\subset E$, at distance $R$ from $S$, is homotopic to a simple closed
curve $\alpha'$ on $S$, and let $C_\alpha$ be the ruled cylinder between them.
If a curve $\beta\subset E$ is disjoint from $C_\alpha$ and is homotopic to a
simple closed curve $\beta'$ on $S$, then
\[
 i(\alpha',\beta')\leq
 K\bigl[\length(\alpha)+(\length(\alpha)+1)
   (\length(\beta)+e^{-R}+\length(\beta'))\bigr].
\]
\end{proposition}

\begin{theorem}[Uniform intersection-length estimate]
\label{thurston-ch09-measured-ending-estimate}
\dcref{9.3.5}\group{thurston.ch09.ending}\source{thurston:page-0253,thurston:page-0254}\notready
 Let $E$ be a geometrically tame end cut off by $S_{[\epsilon,\infty)}$.
For a geodesic lamination $\gamma$ realized in $E$, let $d_\gamma$ be the
minimum distance from $S_{[\epsilon,\infty)}$ to an uncrumpled surface through
$\gamma$.  There is a constant $K$ such that any two realized measured
laminations $(\gamma_1,\mu_1),(\gamma_2,\mu_2)\in\mathcal{ML}_0(S)$ satisfy
\[
 i\bigl((\gamma_1,\mu_1),(\gamma_2,\mu_2)\bigr)
 \leq K e^{-2R}\length_S(\gamma_1,\mu_1)\length_S(\gamma_2,\mu_2),
 \qquad R=\inf\{d_{\gamma_1},d_{\gamma_2}\}.
\]
\end{theorem}
\begin{proof}\sketch Apply the curve estimate first to simple closed geodesics that are not short.  The distance bound compares their lengths in $E$ with their lengths on $S$ and supplies one factor of $e^{-R}$ for each curve.  Homogeneity and continuity of length and intersection number extend the inequality to all of $\mathcal{ML}_0(S)$.\end{proof}

\begin{definition}[Ending lamination]
\label{thurston-ch09-ending-lamination}
\dcref{9.3.6}\group{thurston.ch09.ending}\source{thurston:page-0256}\notready
 The ending lamination $\epsilon(E)$ of a geometrically infinite end $E$ is the union of all limits in $\mathcal{GL}(S)$ of sequences of closed geodesics exiting every compact subset of $E$.
\end{definition}

\begin{proposition}[Structure of the ending lamination]
\label{thurston-ch09-ending-lamination-structure}
\dcref{9.3.7,9.3.8,9.3.9}\group{thurston.ch09.ending}\source{thurston:page-0257}\source{thurston:page-0258}\notready
 For every compact subset of $E$, sufficiently far-out bounded-length curves converge to $\epsilon(E)$. Every leaf of $\epsilon(E)$ is dense in it, and every nontrivial closed curve intersects it. A compactly supported lamination admitting arbitrarily short realizations exiting the end is therefore a sublamination of $\epsilon(E)$.
\end{proposition}

\begin{corollary}[Ending lamination forces degeneration]
\label{thurston-ch09-ending-lamination-corollary}
\dcref{9.3.10}\group{thurston.ch09.ending}\source{thurston:page-0259}\notready
 Any injective homotopy class $f:S\to N$ whose images exit a geometrically infinite end determines the same ending lamination; no compactly supported essential lamination disjoint from it can occur.
\end{corollary}

\section{Taming the topology of an end}

\begin{theorem}[Geometric tameness implies topological tameness]
\label{thurston-ch09-geometric-implies-topological-tameness}
\dcref{9.4.1}\group{thurston.ch09.tameness}\source{thurston:page-0256}\source{thurston:page-0257}\notready
 A geometrically tame end of $N-P$ is homeomorphic, outside a compact set, to $S\times[0,\infty)$.
\end{theorem}
\begin{proof}\sketch Exhaust the end by properly embedded uncrumpled surfaces with controlled geometry; the ruled regions between successive surfaces give product blocks, and the blocks assemble to a product outside a compact set.\end{proof}

\begin{corollary}[Almost tameness implies tameness]
\label{thurston-ch09-almost-implies-tame}
\dcref{9.4.2}\group{thurston.ch09.tameness}\source{thurston:page-0257}\notready
 Every almost geometrically tame end is geometrically tame.
\end{corollary}

\begin{proposition}[Proper surface interpolation]
\label{thurston-ch09-proper-interpolation}
\dcref{9.4.3}\group{thurston.ch09.tameness}\source{thurston:page-0258}\source{thurston:page-0259}\notready
 If a proper map $F:S_{[\epsilon,\infty)}\times[0,\infty)\to E$ has uniformly controlled negatively curved slices and exits every compact set in the end direction, then it can be homotoped to a proper product structure, yielding topological tameness.
\end{proposition}

\section{Interpolating negatively curved surfaces}

\begin{definition}[The lamination locus $\mathcal T_\gamma$]
\label{thurston-ch09-t-gamma}
\group{thurston.ch09.interpolation}\source{thurston:page-0258}\source{thurston:page-0260}\notready
 For $\gamma\in\mathcal{ML}_0(S)$, let $\mathcal T_\gamma$ be the set of limiting geodesic laminations of negatively curved uncrumpled surfaces whose wrinkling or bending data tends to $\gamma$.
\end{definition}

\begin{proposition}[Punctured-torus lamination space]
\label{thurston-ch09-punctured-torus-pl}
\dcref{9.5.2,9.5.3}\group{thurston.ch09.interpolation}\source{thurston:page-0260}\source{thurston:page-0261}\notready
 For the once-punctured torus $T-p$, every projective measured lamination is represented by a slope and $\mathcal{PL}_0(T-p)$ is a circle. The locus $\mathcal T_\gamma$ is determined by the complementary slope data.
\end{proposition}

\begin{proposition}[Interpolation on non-punctured-torus surfaces]
\label{thurston-ch09-negative-surface-interpolation}
\dcref{9.5.4}\group{thurston.ch09.interpolation}\source{thurston:page-0262}\source{thurston:page-0263}\notready
 On any hyperbolic surface other than a punctured torus, laminations in the relevant projective class can be interpolated by negatively curved uncrumpled surfaces; the construction is unique near an essentially complete lamination.
\end{proposition}
\begin{proof}\sketch Cut along the lamination, insert negatively curved ruled pieces in complementary polygons, and solve the cusp angle constraints. The local deformation parameters are train-track weights, so interpolation is piecewise linear.\end{proof}

\begin{proposition}[Cusp angle balance]
\label{thurston-ch09-cusp-angle-balance}
\dcref{9.5.5}\group{thurston.ch09.interpolation}\source{thurston:page-0264}\notready
 The sum of dihedral angles along all edges of the wrinkling locus tending to any cusp is zero in $\R/2\pi\Z$.
\end{proposition}

\begin{corollary}[Uniqueness near cusps]
\label{thurston-ch09-uncrumpled-uniqueness}
\dcref{9.5.6}\group{thurston.ch09.interpolation}\source{thurston:page-0264}\source{thurston:page-0265}\notready
 An uncrumpled surface realizing an essentially complete lamination in a fixed homotopy class is unique and is totally geodesic near its cusps.
\end{corollary}

\begin{proposition}[Stability of uncrumpled realizations]
\label{thurston-ch09-uncrumpled-stability}
\dcref{9.5.7}\group{thurston.ch09.interpolation}\source{thurston:page-0265}\source{thurston:page-0266}\notready
 If an essentially complete lamination $\gamma$ is realized by $U$, then every uncrumpled surface realizing a lamination sufficiently close to $\gamma$ is close to $U$.
\end{proposition}

\begin{proposition}[PL structure on measured laminations]
\label{thurston-ch09-ml-pl-structure}
\dcref{9.5.8}\group{thurston.ch09.interpolation}\source{thurston:page-0266}\notready
 The spaces $\mathcal{ML}(S)$ and $\mathcal{ML}_0(S)$ carry canonical piecewise-linear structures supplied by train-track coordinates.
\end{proposition}

\section{Realizations for surface groups with extra parabolics}

\begin{proposition}[Realizability criterion]
\label{thurston-ch09-realizability-criterion}
\dcref{9.7.1}\group{thurston.ch09.realizations}\source{thurston:page-0277}\source{thurston:page-0278}\notready
 For a surface group representation with extra parabolics, a lamination $\gamma\in\mathcal{GL}_0(S)$ is realizable in $N$ exactly when it does not contain a component of either obstruction lamination $\chi_+$ or $\chi_-$ associated to the two ends.
\end{proposition}

\begin{proposition}[Train-track approximation]
\label{thurston-ch09-train-track-approximation}
\dcref{9.7.2}\group{thurston.ch09.realizations}\source{thurston:page-0278}\source{thurston:page-0279}\notready
 Every essentially complete lamination admits a sequence of train tracks carrying it, with complementary regions and transverse weights converging to the lamination; the associated coordinate charts exhaust the relevant projective lamination space.
\end{proposition}

\begin{corollary}[Projective lamination sphere]
\label{thurston-ch09-projective-lamination-sphere}
\dcref{9.7.3}\group{thurston.ch09.realizations}\source{thurston:page-0279}\notready
 $\mathcal{PL}_0(S)$ is homeomorphic to a sphere.
\end{corollary}

\begin{proposition}[Intersection formula in train-track coordinates]
\label{thurston-ch09-train-track-intersection}
\dcref{9.7.4}\group{thurston.ch09.realizations}\source{thurston:page-0285}\source{thurston:page-0286}\notready
 Let $\tau$ be an essentially complete train track.  If
$\gamma\in V_\tau$ and $\gamma^*\in V_\tau^*$ have respectively transverse
branch weights $\mu_\gamma$ and tangential branch weights $\nu_{\gamma^*}$,
then
\[
 i(\gamma,\gamma^*)=
 \sum_{b\text{ a branch of }\tau}\mu_\gamma(b)\nu_{\gamma^*}(b).
\]
Consequently intersection number is bilinear on
$V_\tau\times V_\tau^*\subset\mathcal{ML}_0(S)\times\mathcal{ML}_0(S)$.
\end{proposition}

\begin{corollary}[Local bilinearity near nonconjugate laminations]
\label{thurston-ch09-nonconjugate-lamination-separation}
\dcref{9.7.5}\group{thurston.ch09.realizations}\source{thurston:page-0286}\notready
 If $\gamma,\gamma'\in\mathcal{ML}_0(S)$ are not topologically conjugate and
at least one is essentially complete, then there are neighborhoods $U$ of
$\gamma$ and $U'$ of $\gamma'$, each equipped with linear coordinates, for
which intersection number restricts to a bilinear function on $U\times U'$.
\end{corollary}

\begin{proposition}[Transverse image of $S_\gamma$]
\label{thurston-ch09-transverse-lamination-image}
\dcref{9.7.6}\group{thurston.ch09.realizations}\source{thurston:page-0286,thurston:page-0287}\notready
 Let $\gamma$ be essentially complete, let
$\tau_1\to\tau_2\to\cdots$ define its stereographic coordinate system
$S_\gamma$, and let $\Delta_\gamma$ be the projective classes of measures
supported on $\gamma$.  For each $\gamma'\in\Delta_\gamma$, choose a
tangential measure $\nu_{\gamma'}$ on $\tau_1$ representing $\gamma'$.  The
image of $S_\gamma$ consists exactly of the transverse, possibly nonpositive,
measures $\mu$ on $\tau_1$ satisfying
\[
 \nu_{\gamma'}\mathbin{\cdot}\mu>0
 \quad\text{for every }\gamma'\in\Delta_\gamma.
\]
\end{proposition}

\begin{proposition}[Convexity of unrealizable regions]
\label{thurston-ch09-unrealizable-convexity}
\group{thurston.ch09.realizations}\source{thurston:page-0288}\source{thurston:page-0289}\source{thurston:page-0290}\source{thurston:page-0291}\notready
 For a connected lamination $\gamma$, the set $Z_\gamma$ of laminations intersecting $\gamma$ in the prescribed unrealizable manner is convex in each train-track or stereographic chart. For nonconnected $\gamma$ it is the intersection of the corresponding connected sets. Hence each obstruction set $U_\pm$, formed from the finitely many components of $\chi_\pm$, is a finite union of convex pieces of codimension at least one.
\end{proposition}
\begin{proof}\sketch For a simple curve, corridor trajectories give linear inequalities. For a connected nonsimple lamination, write $Z_\gamma$ as the convex join of measures on $\gamma$ with laminations supported in its complementary subsurface; intersect componentwise in the disconnected case.\end{proof}

\section{Ergodicity of the geodesic flow}

\begin{theorem}[Sullivan's recurrence criterion]
\label{thurston-ch09-sullivan-ergodicity}
\dcref{9.9.1}\group{thurston.ch09.ergodicity}\source{thurston:page-0293}\source{thurston:page-0294}\source{thurston:page-0295}\source{thurston:page-0296}\source{thurston:page-0297}\source{thurston:page-0298}\notready
 Let $M^n$ be complete hyperbolic. The following are equivalent: (a) $\sum_{\gamma\in\pi_1M}\exp(-(n-1)d(x_0,\gamma x_0))$ diverges; (b) the geodesic flow is not dissipative; (c) the geodesic flow on $T_1M$ is recurrent; (d) that flow is ergodic.
\end{theorem}
\begin{proof}\sketch (d)$\Rightarrow$(c) is the general recurrence implication, and (c)$\Rightarrow$(b) is immediate. Relate divergence of the series to infinite visual area of translates of a ball: recurrent rays imply divergence, while dissipativity permits a maximal separated family of translates whose visual shadows decrease geometrically, proving convergence. Finally induce the flow on a positive-measure recurrent return set and use the north-south behavior of large group elements on pairs of ideal endpoints; density points force every invariant measurable set to have measure zero or full measure.\end{proof}

\begin{corollary}[Superharmonic function from nonergodicity]
\label{thurston-ch09-superharmonic-function}
\dcref{9.9.2}\group{thurston.ch09.ergodicity}\source{thurston:page-0294}\notready
 If the geodesic flow is not ergodic, $M$ admits a nonconstant bounded superharmonic function.
\end{corollary}
\begin{proof}\sketch Convergence of the critical Green-function series produces a Green function $G$ on $M$; composing with the convex bounded function $\arctan$ gives a bounded superharmonic function.\end{proof}

\begin{corollary}[Ergodicity for geometrically tame groups]
\label{thurston-ch09-tame-ergodicity-limit-set}
\dcref{9.9.3}\group{thurston.ch09.ergodicity}\source{thurston:page-0294}\notready
 If $\Gamma$ is geometrically tame, the geodesic flow on $T_1(H^n/\Gamma)$ is ergodic exactly when its limit set is all of $S^2$.
\end{corollary}
